\documentclass{article}
\usepackage{blindtext}
\usepackage{multicol}
\title{Útvonalfeldolgozási megoldások napjainkban}
\author{Csuport Ákos \\ email {ckcbs0@inf.elte.hu}}
\affil[1]{Eötvös Loránd Tudományegyetem Informatikai Kar}
\date{2026. 10. 07.}

\begin{document}
\maketitle

\begin{multicols}{2}
[
\section{Bevezetés}
]
Az útvonalak elemzése az informatikában egy különösen aktív terület. Számos algoritmus és módszer jelent meg megtett útvonalak
különböző fajtájú feldolgozására: összegzés, klasszifikáció, predikció, stb. \cite{hussein2026kafy}

Indokolhatja a terület felkapottságát, hogy számos területen, főleg olyanokon, amelyek valamilyen módon a közlekedési infrastruktúrát érintik,
szükség van útvonaladatok összegzésére, összehasonlítására, következő lépések megjóslására, \cite{fang2021mdtp} vagy tervezés alatt lévő megoldások esetén előzetes forgalmi modellek 
létrehozására a még nem kivitelezett hálózatokon \cite{hussein2026kafy}. Ilyenek többek közt a településfejlesztés \cite{zhou2024red} \cite{hussein2026kafy}, útvonaltervezési szolgáltatások \cite{hussein2026kafy} vagy a közlekedésszervezés 
optimalizációja \cite{zhou2024red}.

Eleinte az utak feldolgozására gyakorta használtak olyan tradicionális gépi tanulási technikákat, mint a probabilisztikus modellek, 
bayes-i hálózatok, vagy a \textit{k} legközelebbi szomszéd algoritmusa \cite{fang2021mdtp}. A neurális hálók megjelenését követően ezek legkülönbözőbb változatait
lett szokásos felhasználni útvonalelemzésekhez \cite{hussein2026kafy} \cite{zhou2024red} \cite{fang2021mdtp}.

Az útvonalak, illetve szegmenseinek vektorokba történő beágyazását, és neuronhálókkal történő feldolgozását útvonalreprezentációs
tanulásnak (\textit{TRL, Trajectory Representation Learning}) szokás hívni \cite{zhou2024red}. Az ilyen megoldások mögött többnyire rekurrens neuronhálók
állnak, de egyes kivitelezésekben a transzformer modellt is alkalmazzák. \cite{zhou2024red} Ez a megoldás népszerű és hatékony bizonyos útvonalakkal kapcsolatos
elemzések végrehajtására, de a TRL megoldások jellemzően csak bizonyos feladattípusok megoldására jók (hasonlóság, útvonalak klasszifikációja,
útvonalak úthálózatra igazítása) \cite{hussein2026kafy} \cite{zhou2024red}, így ezen a megközelítésen kívül is bőven van helye útvonalelemzési technikáknak. Ebbe beletartoznak
gráfok feldolgozására szolgáló konvolúciós hálók \cite{fang2021mdtp}, illetve a TRL-nél már említett transzformer modell \cite{vaswani2017attention} erős és utóbbi időben nagy nyelvi
modellekhez gyakran alkalmazott természetes nyelvfeldolgozásra való rátermettségének felhasználása az útvonalak nyelvként történő formalizálásával,
ami bemenetként használható.

\section{Kapcsolódó művek}

Mivel az útvonalak feldolgozása egy kellőképp aktív téma, számos cikk foglalkozott vele akár csak az utóbbi években is \cite{hussein2026kafy}. A RED \cite{zhou2024red}
(\textit{Road-aware masking strategy, Embedding, Dual-objective task learning}) egy önfelügyelő tanulást alkalmazó, transzformer modellre
épülő TRL-keretrendszer. Mint a többi TRL megközelítés, az útvonalak egyes szegmenseit vektorokba ágyazza, az úthálózatra illeszti a
megtett útról gyűjtött helyadatokat, és két feladattal tanítja magát: a RED transzformere az útvonal következő szegmenseit próbálja előre eltalálni, 
a dekóder pedig a teljes útvonalakat rekonstruálja. \cite{zhou2024red}

A betanított modell számos feladatra képes, amit egy TRL megoldás nyújthat. Az útvonalreprezentációs tanulás lehetőségeit kihasználva
az útvonalak olyan szempontjai szerint is képes hasonlítást, klasszifikációt és predikciót végezni, melyet legfeljebb egy-két
vetélytárs tud a kategóriában. Ebbe beletartozik az útvonalak nemcsak térbeli, hanem időbeli elemzése is, valamint az útvonalak felhasználókhoz is társítva
vannak, ezáltal felhasználói preferenciákat is képes kimutatni a keretrendszer. Az eredmények alapján a RED a legtöbb esetben több mint 5%-kal 
javítja a legjobb baseline módszer által elért pontosságot \cite{zhou2024red}.

Éppen a TRL modellek korlátozásainak átlépésére készült egy másik technika, a KAFY \cite{hussein2026kafy}. A KAFY a RED-hez hasonlít annyiban, hogy időben és térben is képes útvonalakat feldolgozni,
ráadásul szintén felhasználja a transzformer modellt, viszont az útvonal szegmenseinek vektorokba ágyazása helyett az útvonalakat egy nyelvként \cite{devlin2018bert} kódolja, majd (emberi nyelven
még nem betanított) transzformereket használ fel. Ennek célja a transzformer modell teljesítményének kihasználása a nyelvfeldolgozásban, amit olyan rendszerek, mint a BERT \cite{devlin2018bert}
vagy a GPT \cite{yenduri2023generative} korábban már demonstráltak. Maga a kialakított nyelv az útvonalak és emberi nyelveken kifejezett szövegek közti
megfeleltetésekre épül: a nyelv különböző nyelvi elemek (pl. szavak) sorozata, az útvonal pedig koordináták sorozata. Egy szövegben,
hogy értelmezhessük, a nyelvi elemek kapcsolatainak szemantikailag helyesnek kell lenniük, hasonlóképp van az útvonal egyes pontjain
készült tér- és időbeli pontok szemantikát alkotnak. Így, míg a TRL rendszereknél előre meghatározott, legtöbbször egy vagy néhány,
feladatra betanított neuronhálók vannak a modellhez szorosan kötődő vektorbeágyazással, a KAFY egy általános, NLP-ben használatoshoz
hasonló tokensorozatot készít az útvonalakból, amit a KAFY-ban demonstrált három transzformer (BERT, GPT-2 és T-5) öt feladatot
lát el: összegzés, predikció, útvonal-generálás, imputáció és klasszifikáció \cite{hussein2026kafy}, azonban a KAFY moduláris és kiegészíthető további, egyedileg
hangolt transzformerrel, melyek általánosan fel tudják dolgozni az útvonalat reprezentáló tokensorozatot.

Egy harmadik útvonalfeldolgozási modell, az MDTP (\textit{Multi-source Deep Traffic Prediction Framework}) \cite{fang2021mdtp} az eddigi említett
megközelítésekkel ellentétben a vizsgált területet régiókra bontja, amik gráfot alkotnak, és a köztük haladó éleken van reprezentálva
a köztük haladó forgalom fajtája (pl. taxi, kerékpár) és mennyisége, időintervallumokra leosztva. Az MDTP korrelációkat mutat ki időszak
és járműtípus szempontjából a különböző régiók közti forgalomban, erre tartalmaz egy komponenst, mely vizualizálja az adatokat,
illetve az eltelt időintervallumok alapján becsli a következő intervallumban lebonyolodó forgalmat \cite{fang2021mdtp}. Ehhez egy gráfkonvolúciós hálót
használ (GCN), ami jó választás egy térben rögzített területen a forgalom régiók közti lebonyolódásának elemzésére \cite{fang2021mdtp}.

\end{multicols}

\end{document}

[1] - RED
[2] - KAFY
[3] - MDTP
[4] - Transformer
[5] - BERT
[6] - GPT